\chapter{Snoring: rumore nei dataset di defect prediction}

\section{Introduzione e contesto}

La \textbf{defect prediction} mira a identificare gli artefatti software, come classi o file, che hanno un'alta probabilità di presentare un difetto. Il suo obiettivo principale è ridurre i costi di testing e code review, permettendo agli sviluppatori di concentrare gli sforzi sulle parti più rischiose del codice.

Il contesto dello studio è la \textbf{within-project defect prediction}, cioè la classificazione delle classi difettose all'interno di specifiche release di uno stesso progetto.

L'affidabilità di un modello predittivo dipende direttamente dalla qualità del dataset. Se il dataset contiene rumore, anche la valutazione e l'addestramento dei classificatori possono risultare distorti.

\begin{notaprof}
Il concetto centrale è il principio \textbf{Garbage In, Garbage Out}: se le etichette del dataset sono sbagliate, anche il modello più sofisticato apprenderà informazioni distorte.
\end{notaprof}

\begin{notaslide}
Le slide ricordano che lavori precedenti avevano già individuato altre fonti di rumore nei dataset di defect prediction, come la misclassification e l'origine errata dei difetti.
\end{notaslide}

\section{Sleeping defects e snoring classes}

Uno \textbf{sleeping defect}, o \textbf{dormant bug}, è un difetto che viene scoperto o corretto diverse release dopo la sua introduzione effettiva nel codice. Il difetto è quindi già presente nel sistema, ma non è ancora visibile al costruttore del dataset.

Il ciclo di vita di un difetto può essere descritto tramite tre eventi principali:

\begin{itemize}
    \item \textbf{I, Injected}: il difetto viene introdotto nel codice;
    \item \textbf{N oppure OV, Noticing/Opening Version}: il difetto viene osservato o viene aperto il ticket;
    \item \textbf{F oppure FV, Fixed Version}: il difetto viene corretto tramite fix commit.
\end{itemize}

Le \textbf{snoring classes} sono classi affette esclusivamente da difetti non ancora corretti. Poiché l'esistenza del difetto non è conoscibile prima del fix, queste classi vengono erroneamente considerate non difettose.

\begin{notaprof}
La metafora è utile: il difetto \emph{dorme}, mentre la classe \emph{russa}, perché produce rumore nel dataset. Una classe russa solo se tutti i difetti che la colpiscono sono ancora dormienti. Se contiene almeno un difetto già noto e corretto, viene comunque etichettata come difettosa.
\end{notaprof}

\begin{notaesame}
Lo \textbf{snoring} introduce nei dataset \textbf{falsi negativi}, non falsi positivi. Una classe etichettata come difettosa è supportata da un fix osservato; una classe etichettata come non difettosa, invece, potrebbe contenere un difetto non ancora scoperto.
\end{notaesame}

\begin{notaslide}
Le slide mostrano un esempio con tre classi, C1, C2 e C3, osservate su più release. Una classe può essere realmente difettosa in una release intermedia, ma risultare \emph{not defective} se il relativo fix avviene solo in una release successiva.
\end{notaslide}

\section{Assegnazione dei difetti alle classi}

Per stabilire da quando una classe deve essere considerata difettosa, il processo può usare due fonti principali:

\begin{enumerate}
    \item la prima \textbf{Affected Version} riportata nel sistema di issue tracking, ad esempio Jira;
    \item un algoritmo di ricostruzione storica, come \textbf{SZZ}, quando l'Affected Version non è disponibile o non è affidabile.
\end{enumerate}

Le nozioni principali sono:

\begin{itemize}
    \item \textbf{IV, Injected Version}: release in cui il difetto viene introdotto;
    \item \textbf{OV, Opening Version}: release in cui viene aperto il ticket;
    \item \textbf{FV, Fixed Version}: release in cui il difetto viene corretto;
    \item \textbf{AV, Affected Version}: informazione che indica le versioni affette dal difetto.
\end{itemize}

\begin{notaprof}
Una classe può essere realmente difettosa ma apparire non buggy se, nello snapshot analizzato, il difetto è già presente ma il relativo ticket o fix sarà osservato solo in futuro. Questo è il punto in cui nasce il falso negativo.
\end{notaprof}

\section{SZZ: funzionamento e limiti}

\textbf{SZZ} è un processo algoritmico usato per identificare quando un difetto è stato introdotto in un progetto software.

A partire da un \textbf{bug-fix commit}, SZZ analizza le linee modificate dalla correzione e usa meccanismi come \texttt{git blame} per risalire all'ultima modifica precedente di quelle stesse linee. Tale modifica viene interpretata come possibile \textbf{bug-introducing change}.

\begin{notaslide}
Nelle slide è mostrato un esempio in cui un bug di tipo ``Array index out of range'' viene corretto modificando la condizione di un ciclo da \texttt{i<=4} a \texttt{i<4}. SZZ parte dal fix e risale alla versione precedente in cui la riga errata era stata introdotta.
\end{notaslide}

SZZ presenta diversi limiti:

\begin{itemize}
    \item può fallire quando il fix consiste in una \textbf{code addition};
    \item è sensibile ai \textbf{refactoring};
    \item può produrre una \textbf{backward search} imprecisa;
    \item non gestisce perfettamente i \textbf{regressive bugs};
    \item non elimina il problema dello \textbf{snoring}, perché i difetti non ancora corretti restano invisibili.
\end{itemize}

\begin{notaslide}
Le slide riportano che solo una parte dei bug-introducing changes può essere individuata tramite l'assunzione classica di SZZ, secondo cui il bug è stato introdotto nelle stesse linee poi modificate dal fix.
\end{notaslide}

\section{Research Questions dello studio}

Lo studio sullo snoring si articola in cinque research questions:

\begin{enumerate}
    \item \textbf{RQ1: To what extent do defects sleep?}
    \item \textbf{RQ2: To what extent do classes snore?}
    \item \textbf{RQ3: To what extent does snoring impact the evaluation of classifiers accuracy?}
    \item \textbf{RQ4: To what extent does snoring impact defect prediction accuracy?}
    \item \textbf{RQ5: To what extent is no data better than snoring data?}
\end{enumerate}

Le prime due domande misurano l'entità del fenomeno, mentre le ultime tre analizzano l'impatto dello snoring sulla valutazione e sull'addestramento dei classificatori.

\section{RQ1: To what extent do defects sleep?}

\subsection{Rationale}

La prima research question misura quante release trascorrono tra l'introduzione di un difetto e la sua correzione. L'obiettivo è quantificare quanto a lungo un difetto possa restare invisibile nel progetto.

\subsection{Design}

Sono stati selezionati \textbf{20 grandi progetti Apache}, con issue tracking in Jira e codice versionato in Git. I progetti considerati hanno almeno 10 release e un buon collegamento tra bug e codice.

\subsection{Variabili}

\begin{itemize}
    \item \textbf{Variabile indipendente}: il bug o ticket in Jira.
    \item \textbf{Variabile dipendente}: numero di release per cui il bug ha dormito, misurato tra l'introduzione e il fix.
\end{itemize}

\subsection{Risultati e discussione}

Il risultato principale è che il bug medio dorme per circa \textbf{7 release}. Inoltre, nella maggior parte dei progetti, i difetti dormono per oltre il \textbf{19\%} delle release esistenti.

\begin{notaslide}
Le slide riportano casi estremi come \textbf{LUCENE-3672}, dormiente per 84 release, e \textbf{STRATOS-1653}, dormiente per 51 release su 53.
\end{notaslide}

\section{RQ2: To what extent do classes snore?}

\subsection{Rationale}

Non ogni sleeping defect produce necessariamente una snoring class. Se una classe contiene anche un altro difetto già noto, la classe viene comunque etichettata come difettosa e non genera falso negativo. La RQ2 misura quindi quanto i difetti dormienti si traducano effettivamente in classi rumorose.

\subsection{Design}

Lo studio osserva come cambia la difettosità stimata man mano che l'osservatore avanza nel tempo. Nelle slide questo è rappresentato con i \textbf{binocoli}: più ci si sposta avanti nella storia del progetto, più difetti prima invisibili diventano osservabili.

\subsection{Variabili}

\begin{itemize}
    \item \textbf{Variabile indipendente}: \textbf{Progress}, cioè il rapporto tra la release osservata e il numero totale di release.
    \item \textbf{Variabile dipendente}: \textbf{Missing Rate}, cioè:
    \[
        \frac{FN}{FN + TP}
    \]
\end{itemize}

Il Missing Rate rappresenta la quota di classi realmente difettose che non vengono riconosciute come tali. Corrisponde a \(1 - Recall\).

\subsection{Risultati e discussione}

Il Missing Rate diminuisce con il progresso temporale, ma resta significativo per molte release. Per la maggior parte dei progetti, il Missing Rate rimane superiore al \textbf{50\%} se non si rimuove più del \textbf{20\%} delle release finali.

\begin{notaprof}
Per ottenere un Missing Rate nullo bisogna in media ignorare una porzione molto ampia delle release finali. Questo mostra che la defectiveness osservata si stabilizza solo molto tempo dopo l'introduzione reale dei difetti.
\end{notaprof}

\section{RQ3: Impatto sulla valutazione dei classificatori}

\subsection{Rationale}

La terza research question analizza se lo snoring alteri la valutazione dei classificatori. Se i dati di test sono rumorosi, si rischia di scegliere come migliore un classificatore che in realtà non lo è.

\subsection{Design}

Sono stati valutati \textbf{15 classificatori} tramite \textbf{ordered holdout}, mantenendo l'ordine cronologico dei dati: 66\% per il training e 33\% per il testing.

\subsection{Metriche}

Le metriche considerate includono:

\begin{itemize}
    \item Precision;
    \item Recall;
    \item F1;
    \item Cohen's Kappa;
    \item AUC;
    \item Matthews Correlation Coefficient.
\end{itemize}

Viene inoltre calcolato il \textbf{Relative Bias}, cioè la distanza relativa tra accuratezza misurata su dati rumorosi e accuratezza misurata su dati privi di snoring.

\subsection{Risultati e discussione}

Lo snoring introduce un bias elevato nella valutazione. In particolare, usando dati affetti da snoring, il classificatore realmente migliore viene selezionato correttamente solo in una parte dei casi.

\begin{notaslide}
Le slide riportano che il classificatore migliore viene selezionato correttamente solo nel \textbf{45\%} dei casi quando la valutazione è eseguita su dati affetti da snoring.
\end{notaslide}

\section{RQ4: Impatto sull'accuratezza predittiva}

\subsection{Rationale}

La quarta research question valuta quanto lo snoring peggiori l'accuratezza predittiva reale dei modelli.

\subsection{Design}

I modelli vengono addestrati su dati affetti da snoring e poi valutati su un test set pulito. In questo modo si misura l'effetto del rumore presente nel training set.

\subsection{Risultati e discussione}

Lo snoring degrada le performance di tutti i classificatori e di tutte le metriche osservate. Le metriche più colpite sono \textbf{Recall}, \textbf{F1}, \textbf{Kappa} e \textbf{Matthews}. La \textbf{Precision} risulta meno sensibile.

\begin{notaesame}
Lo snoring danneggia soprattutto la \textbf{Recall}, perché introduce falsi negativi nel training set. Il modello impara quindi a classificare come sane molte classi che in realtà sono difettose.
\end{notaesame}

\section{RQ5: Quando no data è meglio di snoring data}

\subsection{Rationale}

La quinta research question valuta se sia preferibile rimuovere alcune release recenti per ridurre il rumore, anche al costo di avere meno dati.

\subsection{Design}

Lo studio confronta scenari in cui vengono rimosse le ultime \textbf{1, 2, 3 o 4 release} dal training set.

\subsection{Risultati e discussione}

Rimuovere una release tende a migliorare la qualità della prediction. Tuttavia, rimuovere troppe release peggiora le performance perché riduce eccessivamente la dimensione del dataset.

\begin{notaprof}
Il risultato mostra un trade-off tra \textbf{Garbage In, Garbage Out} e \textbf{dimensione del dataset}. Pochi dati puliti possono essere meglio di molti dati rumorosi, ma eliminare troppe release fa perdere informazione utile.
\end{notaprof}

\section{Conclusioni}

Lo snoring è una fonte rilevante di rumore nei dataset di defect prediction. Il fenomeno nasce dal fatto che molti difetti restano dormienti per diverse release e diventano osservabili solo dopo il fix.

I risultati principali sono:

\begin{itemize}
    \item molti difetti dormono per una porzione significativa della vita del progetto;
    \item le snoring classes introducono falsi negativi;
    \item la valutazione dei classificatori può diventare fortemente biasata;
    \item la Recall è particolarmente danneggiata;
    \item rimuovere alcune release recenti può migliorare la qualità del training set, ma rimuoverne troppe peggiora il modello.
\end{itemize}

\section{Future works}

Le possibili estensioni includono:

\begin{itemize}
    \item tecniche più fine-grained per identificare e rimuovere il rumore;
    \item combinazione di più strategie di mitigazione;
    \item uso di etichette probabilistiche invece di etichette binarie rigide;
    \item replica dello studio nel contesto della \textbf{Just-In-Time defect prediction}.
\end{itemize}

\section{Da ricordare}

\begin{notaesame}
Lo snoring introduce falsi negativi: classi realmente difettose possono essere etichettate come non difettose perché il fix avverrà solo in futuro.
\end{notaesame}

\begin{notaesame}
La metrica più penalizzata è la Recall, perché il modello impara a non riconoscere molti veri positivi.
\end{notaesame}

\begin{notaesame}
Rimuovere alcune release recenti può ridurre il rumore, ma eliminare troppe release riduce eccessivamente la quantità di dati disponibili.
\end{notaesame}